\textbf{Generalization to completely held-out prescribed-fire families.}
\textbf{a}, Mean absolute error (MAE) on each held-out family for all seven prespecified models, averaged over five independently trained seeds. Family sample counts are shown in parentheses. The final two rows are descriptive averages over the ten separately trained holdout experiments: the case-weighted average assigns each of the 997 cross-held-out cases equal weight, whereas the equal-family average assigns each of the ten families equal weight. They do not represent predictions from one common trained checkpoint, and no aggregate confidence interval is asserted.
\textbf{b,c}, The two prespecified paired MAE contrasts for every family. Diamonds show the mean paired effect; lines show the existing percentile 95\% intervals from 5,000 resamples of whole geometry clusters and training seeds. Positive differences favor the second named representation. All intervals are shown on a common scale, without multiplicity adjustment; they do not establish a global family-population claim. Individual time points are not resampled.
All errors and differences are percentage points of the dimensionless archived scalar-envelope target. The smoldering holdout retains 86 restored constant-response cases and 11 retained nonconstant cases, all of which contribute to its reported MAE. Its distinct response composition should be considered when interpreting both aggregate rows. The figure includes all 350 completed family-holdout runs and concerns new prescribed boundary histories, not independent physical-capacity validation.
